This project evaluated whether student-accessible quantum computing platforms are currently well enough to implement Grover’s algorithm meaningfully on game-like problems, and at what scale this is feasible. The results show that Grover is currently most useful as a proof-of-concept: small instances can demonstrate correct behavior, but scalability is limited by circuit depth and noise.
\\
\\
For Sudoku, Grover works well for a $2\times 2$ instance on real hardware, but the partially filled $4\times 4$ case remains practical only in simulation and is outperformed by classical structure-aware solvers. For Lights Out, correct amplification is observed in simulation, yet hardware performance degrades strongly due to transpilation overhead and two-qubit gate noise. For WFC, the $3\times 3$ case achieves high success probability in simulation (99.6\%), while hardware results are near-uniform with only $\approx 1.4\%$ valid outcomes, reflecting circuit depths far beyond coherence limits.
\\
\\
Overall, Grover helps when the oracle is compact and the instance is small enough that depth stays within hardware limits. It does not yet provide a practical advantage over classical methods for these structured problems on current student-accessible hardware. Future work should focus on reducing oracle depth, exploring hybrid classical--quantum pipelines (e.g., classical pruning followed by Grover), and repeating these experiments as hardware improves, especially with better connectivity and lower two-qubit error rates.